% ECONOMICS_PROBLEM: {"id": "D44-108", "title": "Simple combinatorial bidding", "jel_code": "D44", "source_id": "OP-0294"}
% !TeX program = xelatex
\documentclass[11pt,a4paper]{article}
\usepackage[margin=23mm]{geometry}
\usepackage{fontspec}
\setmainfont{Latin Modern Roman}
\usepackage{amsmath,amssymb,mathtools}
\usepackage{xcolor,enumitem,booktabs,longtable,array,xurl,hyperref}
\definecolor{muted}{HTML}{53636C}
\hypersetup{colorlinks=true,urlcolor=blue,linkcolor=blue}
\setlength{\parindent}{0pt}
\setlength{\parskip}{5pt}
\newcommand{\R}{\mathbb R}
\newcommand{\E}{\mathbb E}
\newcommand{\N}{\mathbb N}
\newcommand{\ind}{\mathbf 1}
\newcommand{\Var}{\operatorname{Var}}
\newcommand{\Cov}{\operatorname{Cov}}
\newcommand{\supp}{\operatorname{supp}}
\DeclareMathOperator*{\argmin}{arg\,min}
\DeclareMathOperator*{\argmax}{arg\,max}
\newcommand{\entrymeta}[3]{{\sffamily\small\color{muted}\textbf{\texttt{#1}}\quad|\quad #2\par #3\par}}
\newcommand{\Problemref}[1]{\texttt{#1}}
\newcommand{\JELref}[2]{\texttt{#2} (JEL \texttt{#1})}
\begin{document}
\section*{Simple combinatorial bidding}
\textbf{Problem ID:} \texttt{D44-108}\quad \textbf{JEL:} \texttt{D44}\par
% BEGIN ECONOMICS STATEMENT
\label{problem:OP-0294}
{\sffamily\small\color{muted}Original subject group: Mechanism design and market design\par}
\label{definitions:problem:30007576}
\textbf{Audit classification:} Published research agenda. Roth 2010 §1 and the 2024 auction survey §4.3 supply simplicity agendas. The combined cognitive/communication frontier is editorial; unrestricted real-valued messages make a message-count bound meaningless.
\subsection*{Definitions and precise research target}
\textbf{Model and research target.} Let \(G\) be \(m\) indivisible goods and \(n\) bidders have nonnegative monotone values \(v_i:2^G\to[0,1]\). Let \(\mathcal V\) be a prespecified class of valuation profiles admitting both substitutes and complements, with optimal feasible welfare \(\operatorname{OPT}(v)>0\). A bidding interface \(I\) allows at most \(k\) messages, each encoded in at most \(b\) bits, or a separately specified value-query format and has measured decision time \(C_i(I)\). Together with allocation and payment rule \(M\), it induces a distribution of submitted bids and final allocations \(A_i\) for a target bidder population \(P\). Define the performance frontier
\[
 F(k,c)=\sup_{I,M:\sup_{v\in\mathcal V}E_{K_P(\cdot\mid I,M,v)}[C_i(I)]\leq c}
 \inf_{v\in\mathcal V}\frac{E_P[\sum_i v_i(A_i)]}{\max_{A\text{ feasible}}\sum_i v_i(A_i)},
\]
subject to a specified participation and incentive requirement. Declare whether the behavioral response is a rational equilibrium, a calibrated learning model, or an experimentally measured distribution; these interpretations cannot be interchanged. Characterize useful attainable guarantees and complexity costs, design interfaces approaching the frontier, and validate performance under unfamiliar valuations and bidder experience. For strategic guarantees, state the equilibrium concept and exclude arbitrary nonparticipation; for experimental guarantees, quantify sampling uncertainty.

\emph{Definitions and maintained assumptions (editorial unless a source definition is named).} Feasible allocations are disjoint bundles \(A_i\subseteq G\), and \(\mathrm{OPT}(v)=\max_A\sum_i v_i(A_i)\). To make \(k\) constrain communication, restrict each message to an alphabet of at most \(2^b\) symbols (at most \(b\) bits), or specify a concrete bundle/value query format with fixed precision. A single arbitrary real message can encode an entire finite valuation table, so raw message count is not a meaningful information bound. Rational values with input bit length \(B\), or a value-oracle model, specify computational input. For each \((I,M,v)\), define a behavioral kernel \(K_P(\text{bids}\mid I,M,v)\) for the target population. Time expectation must be conditional on valuation profile; use \(\sup_{v\in\mathcal V}\mathbb E_{K_P}[C_i]\le c\) for a worst-case frontier. Incentive and participation conditions must be consistent with that kernel. The ratio is defined only for \(\mathrm{OPT}(v)>0\). The supremum ranges over implementable interfaces satisfying an announced representation/operation budget.
\subsection*{Variants and assumption changes}
\begin{enumerate}
\item \textbf{Strategic bit budget (editorial).} Require dominant-strategy truthfulness and ex-post IR on a specified substitutes or complements domain, with at most \(kb\) communicated bits. Seek worst-case welfare guarantees and information lower bounds.
\item \textbf{Human interfaces (editorial).} Fix a finite set of implementable interfaces and estimate bidding/time kernels experimentally. Compare held-out welfare under identical valuation and population distributions; this empirical frontier cannot automatically inherit rational-equilibrium guarantees.
\end{enumerate}
\subsection*{References and source audit}
\begin{enumerate}
\item Roth §1 p.3; Palacios-Huerta et al. §4.3 pp.23–25. Supports the stated research area; exact displayed specification is editorial. \url{https://web.stanford.edu/~alroth/papers/NSF%20Grand%20Challenge.SBE%202020.%20Market%20Design.pdf}
\item Roth §1 p.3; Palacios-Huerta et al. §4.3 pp.23–25. Supports the stated research area; exact displayed specification is editorial. Original LSE fetch failed; full manuscript recovered from Harvard author page. \url{https://parkes.seas.harvard.edu/sites/g/files/omnuum8711/files/parkes/files/palacios-huerta.pdf}
\end{enumerate}
\begin{enumerate}\raggedright
\item\hypertarget{definitions:source:30007576:1}{} Alvin E. Roth. 2010. \emph{Market Design: Understanding markets well enough to fix them when they’re broken}. §1 Fundamental questions, pp.3–4.
\url{https://web.stanford.edu/~alroth/papers/NSF%20Grand%20Challenge.SBE%202020.%20Market%20Design.pdf}.
\item\hypertarget{definitions:source:30007576:2}{} Ignacio Palacios-Huerta; David C. Parkes; Richard Steinberg. 2024. \emph{Combinatorial Auctions in Practice}. §4.3, printed pp.23–25, Harvard author manuscript; published JEL62(2),517–553 (2024).
\url{https://parkes.seas.harvard.edu/sites/g/files/omnuum8711/files/parkes/files/palacios-huerta.pdf}.
DOI: \url{https://doi.org/10.1257/jel.20221679}.
\end{enumerate}

\subsection*{Common conventions: Source mathematical conventions}

These conventions apply to each entry unless explicitly replaced.

\textbf{Models and identification.} A primitive model \(m\) specifies all probability laws, feasible choices, information, equilibrium and measurement rules used by its target. The admissible class \(\mathcal M\) is fixed independently of outcomes. If \(P_O(m)\) is its observable probability law and \(\tau(m)\) the target, its identified set at observable law \(P\) is
\[
 \mathcal I_\tau(P)=\{\tau(m):m\in\mathcal M,\ P_O(m)=P\}.
\]
Point identification means that this set is a singleton. An empty set signals incompatibility of data and assumptions. Coordinate bounds need not describe the joint set. Bounds are sharp only if every retained target is attainable or arbitrarily approachable by compatible models. Finite bounds require explicit restrictions on unobserved quantities; unrestricted confounding, hidden wealth, extrapolation or arbitrary equilibrium selection can yield uninformative sets.

\textbf{Causal interventions.} A potential outcome \(Y_i(p)\) is the outcome under a complete policy regime \(p\), including timing, financing, information and market spillovers. The target population and units are fixed before treatment. With interference, use the complete assignment vector or a justified exposure mapping; individual no-interference notation alone cannot identify an economy-wide intervention. A difference of mean potential outcomes does not require the cross-world joint distribution. Conditioning on post-treatment migration, survival, enrollment or employment changes the estimand and needs separate justification.

\textbf{Statistical statements.} An estimator, confidence set, forecast or test specifies a sampling law, model class and loss/coverage criterion. Uniform \((1-\alpha)\)-coverage means \(\inf_{m\in\mathcal M}\Pr_m\{\tau(m)\in C_n\}\geq1-\alpha\), or an explicitly asymptotic version. The whole parameter space trivially has coverage, so informativeness requires a stated width, risk or power objective. A forecasting improvement is relative to a named benchmark, information set and evaluation population.

\textbf{Equilibrium and optimization.} An equilibrium comprises feasible plans, optimality, beliefs where needed, budget constraints and all market/resource clearing. The primitives or a stated benchmark define each correspondence \(\mathcal E(p;m)\). Nonemptiness is assumed or proved, never inferred just from compact policy sets. A selected-equilibrium target uses a declared selection rule; a robust target quantifies over all permitted equilibria. Use suprema or infima unless attainment is established. An \(\varepsilon\)-optimal policy is feasible and within \(\varepsilon\) of the finite optimum value. Report an empty feasible set separately.

\textbf{Welfare and accounting.} Normative utility, interpersonal normalization, social weights, numeraire, discounting and the use of public receipts are primitives. Behavioral utility need not coincide with the welfare criterion. Household welfare includes ownership income and rents once; adding profits again to compensating variation generally double-counts them. A monetary equivalent is defined by an explicit transfer and price convention, with monotonicity and range conditions. Average effects, mechanism contributions, transfers and welfare losses are different targets. Infinite sums require convergence and dynamic feasible plans require terminal or no-Ponzi conditions.

\textbf{Source versions.} A working-paper date, revision date and journal publication year are distinguished. A DOI for a journal version is not automatically the DOI of an earlier manuscript. Locators refer to the stated version and printed pages where specified. A metadata-only check does not verify a claim about a full-text conclusion. Every entry's source subsection narrows the provenance claim accordingly.
% END ECONOMICS STATEMENT
\end{document}
